\documentclass{article}
\usepackage{lmodern}
\usepackage{amsmath}
\usepackage{listings}
\usepackage{fullpage}
\usepackage{parskip}
\usepackage{xcolor}
\lstset{
	basicstyle=\ttfamily,
	columns=fullflexible,
	frame=single,
	breaklines=true,
	postbreak=\mbox{\textcolor{red}{$\hookrightarrow$}\space},
}
\begin{document}
\title{Experiment `p\_6\_7\_reverse\_array\_array' Results}
\author{\today}
\date{}
\maketitle
\textbf{Experiment outcome: }SUCCESS
\\\textbf{Bad responses: }0
\\\textbf{Responses containing}~\texttt{assume}~\textbf{: }0
\\\textbf{Responses containing}~\texttt{expect}~\textbf{: }0
\\\textbf{Resolution attempts: }1
\\\textbf{Hard fails (resolution): }0
\\\textbf{Soft fails (resolution): }0
\\\textbf{Verification attempts: }1
\section*{Problem Specification}
\textbf{Problem name: }p\_6\_7\_reverse\_array\_array
\\\textbf{Natural language statement: }Write a method that reverses the sequence of elements in an array. For example, if you call the method with the array [1, 4, 9, 16, 9, 7, 4, 9, 11] then the array is changed to [11, 9, 4, 7, 9, 16, 9, 4, 1].
\\\textbf{Method signature: }p\_6\_7\_reverse\_array\_array(inputs: array<int>) returns (reversed: array<int>)
\subsection*{Ensures}
\begin{itemize}
\item \texttt{reversed.Length == inputs.Length}
\item \texttt{forall i :: 0 <= i < reversed.Length ==> reversed[i] == inputs[inputs.Length - 1 - i]}
\end{itemize}
\clearpage
\section*{GenAI interactions}
Below you will find all interactions between the `user' (program) and the `assistant' (OpenAI).
\subsection*{Program $\rightarrow$ GenAI}
\begin{lstlisting}
You are given the following task to perform in Dafny:

Write a method that reverses the sequence of elements in an array. For example, if you call the method with the array [1, 4, 9, 16, 9, 7, 4, 9, 11] then the array is changed to [11, 9, 4, 7, 9, 16, 9, 4, 1].

The signature should be:

p_6_7_reverse_array_array(inputs: array<int>) returns (reversed: array<int>)

The method should respect the following contract:

ensures reversed.Length == inputs.Length, ensures forall i :: 0 <= i < reversed.Length ==> reversed[i] == inputs[inputs.Length - 1 - i]

Produce and show only the Dafny body of this method, including the curly braces that surround it. Do not show the signature nor contract. You must not use 'assume' nor 'expect' anywhere in your code.
\end{lstlisting}
\subsection*{GenAI $\rightarrow$ Program}
\textbf{System fingerprint: }null
\\\textbf{ID: }chatcmpl-Dn1hMrYlbsoMdkiL9pOtncF2gKBXq
\\\textbf{Model: }gpt-5.5-2026-04-23
\\\textbf{Created at: }1780576972
\\\textbf{Finish reason: }stop
\\\textbf{Usage: }458 tokens in, and 484 tokens out
\begin{lstlisting}
{
  reversed := new int[inputs.Length];
  var i := 0;
  while i < inputs.Length
    invariant 0 <= i <= inputs.Length
    invariant reversed.Length == inputs.Length
    invariant forall j :: 0 <= j < i ==> reversed[j] == inputs[inputs.Length - 1 - j]
  {
    reversed[i] := inputs[inputs.Length - 1 - i];
    i := i + 1;
  }
}
\end{lstlisting}
\section*{Final Program}
\begin{lstlisting}
method p_6_7_reverse_array_array(inputs: array<int>) returns (reversed: array<int>)
	ensures reversed.Length == inputs.Length
	ensures forall i :: 0 <= i < reversed.Length ==> reversed[i] == inputs[inputs.Length - 1 - i]
{
  reversed := new int[inputs.Length];
  var i := 0;
  while i < inputs.Length
    invariant 0 <= i <= inputs.Length
    invariant reversed.Length == inputs.Length
    invariant forall j :: 0 <= j < i ==> reversed[j] == inputs[inputs.Length - 1 - j]
  {
    reversed[i] := inputs[inputs.Length - 1 - i];
    i := i + 1;
  }
}
\end{lstlisting}
\section*{Total Token Usage}
\textbf{Input tokens: }458
\\\textbf{Output tokens: }484
\\\textbf{Reasoning tokens: }361
\\\textbf{Sum of `total tokens': }942
\section*{Experiment Timings}
\textbf{Overall Experiment} started at 1780576971406, ended at 1780577010819, lasting 39413ms (39.41 seconds)
\\\textbf{Iteration \#1} started at 1780576971406, ended at 1780577010819, lasting 39413ms (39.41 seconds)
\end{document}
